\begin{thebibliography}{99}

\bibitem{Repository2026}
\emph{A polynomial definition of the integers in the rationals},
repository \texttt{elliotglazer/diophantine-z-q}, revision
\texttt{28e87b1b62af0effd954605af7f70b2699dd5672}, accessed
6 November 2026. Manuscript \texttt{quartic\_12504.tex}; package
\texttt{quartics/quartic-12504}.
\url{https://github.com/elliotglazer/diophantine-z-q}.
Working fork: \url{https://github.com/JoseBrox/diophantine-z-q_fork}.

\bibitem{Matiyasevich1970}
Yu. V. Matiyasevich, \emph{The Diophantineness of enumerable sets},
Doklady Akademii Nauk SSSR \textbf{191} (1970), no. 2, 279--282;
English translation, Soviet Mathematics Doklady \textbf{11} (1970),
354--357. \url{https://www.mathnet.ru/eng/dan35274}.

\bibitem{MatiyasevichRobinson1975}
Yu. Matiyasevich and J. Robinson, \emph{Reduction of an arbitrary
Diophantine equation to one in 13 unknowns}, Acta Arithmetica
\textbf{27} (1975), 521--553.
\url{https://doi.org/10.4064/aa-27-1-521-553}.

\bibitem{Jones1982}
J. P. Jones, \emph{Universal Diophantine equation}, Journal of Symbolic
Logic \textbf{47} (1982), no. 3, 549--571.
\url{https://doi.org/10.2307/2273588}.

\bibitem{Daans2024}
N. Daans, \emph{Universally defining $\Z$ in $\Q$ with 10 quantifiers},
Journal of the London Mathematical Society \textbf{109} (2024), no. 2,
e12864. \url{https://doi.org/10.1112/jlms.12864}.
Primary preprint: \url{https://arxiv.org/abs/2301.02107}.

\bibitem{DDF2025}
N. Daans, P. Dittmann, and A. Fehm, \emph{Existential rank and
essential dimension of diophantine sets}, arXiv:2102.06941v5 (2025),
in particular Theorem 1.4, Corollary 4.21, and Remark 8.23.
\url{https://arxiv.org/abs/2102.06941v5}.

\bibitem{ZhangSun2022}
G.-R. Zhang and Z.-W. Sun, \emph{$\Q\setminus\Z$ is Diophantine
over $\Q$ with 32 unknowns}, Bulletin of the Polish Academy of Sciences,
Mathematics \textbf{70} (2022), no. 2, 93--106.
\url{https://arxiv.org/abs/2104.02520}.

\bibitem{Sun2026Rational86}
Z.-W. Sun, \emph{86 unknowns for Hilbert's tenth problem over the
rationals}, arXiv:2610.20691v2 (2026).
\url{https://arxiv.org/html/2610.20691v2}.
The new shared-point construction additionally assumes Theorem 4.1,
not merely the algebraic devices of Sections 10.3--10.6. Its use is
explicitly separated from source-preserving algebraic transformations.

\bibitem{FreiTinkova2025}
C. Frei and M. Tinkov\'a, with an appendix by C. Frei, J. K\"onig,
and M. Tinkov\'a, \emph{Linear constellations in primes with arithmetic
restrictions}, arXiv:2408.17441v3 (2025), Theorem 1.4.
\url{https://arxiv.org/html/2408.17441v3}.
The Artin-symbol theorem cited here is unconditional; the separate
primitive-root theorem has a Riemann-hypothesis assumption.

\bibitem{TaoZiegler2018}
T. Tao and T. Ziegler, \emph{Polynomial patterns in the primes},
Forum of Mathematics, Pi \textbf{6} (2018), e1.
\url{https://arxiv.org/abs/1603.07817}.

\bibitem{PatersonStockmeyer1973}
M. S. Paterson and L. J. Stockmeyer, \emph{On the number of nonscalar
multiplications necessary to evaluate polynomials}, SIAM Journal on
Computing \textbf{2} (1973), no. 1, 60--66.
\url{https://doi.org/10.1137/0202007}.

\bibitem{Sun2021}
Z.-W. Sun, \emph{Further results on Hilbert's Tenth Problem},
Science China Mathematics \textbf{64} (2021), no. 2, 281--306;
in particular Theorem 1.1(i) and Equation (1.1).
\url{https://arxiv.org/abs/1704.03504v7}.

\bibitem{Koenigsmann2016}
J. Koenigsmann, \emph{Defining $\Z$ in $\Q$}, Annals of Mathematics
\textbf{183} (2016), no. 1, 73--93.
\url{https://doi.org/10.4007/annals.2016.183.1.2}.
Primary preprint: \url{https://arxiv.org/abs/1011.3424}.

\bibitem{Voight2021}
J. Voight, \emph{Quaternion algebras}, Graduate Texts in Mathematics
\textbf{288}, Springer, 2021; in particular Proposition 14.6.7.
\url{https://doi.org/10.1007/978-3-030-56694-4}.

\bibitem{lbKollar2002}
J. Koll\'ar, \emph{Unirationality of cubic hypersurfaces}, Journal
of the Institute of Mathematics of Jussieu \textbf{1} (2002), no. 3,
467--476, Theorem 2.
\url{https://doi.org/10.1017/S1474748002000117}.
\url{https://arxiv.org/abs/math/0005146}.

\bibitem{lbSchinzel2000}
A. Schinzel, \emph{Polynomials with special regard to reducibility},
Encyclopedia of Mathematics and its Applications \textbf{77},
Cambridge University Press, 2000, Chapter 4, Section 4.
\url{https://doi.org/10.1017/CBO9780511542916}.

\bibitem{lbMuller2001}
P. M\"uller, \emph{Finiteness results for Hilbert's irreducibility
theorem}, arXiv:math/0109071, 2001; introduction and Section 2.1.
\url{https://arxiv.org/abs/math/0109071}.

\bibitem{lbLang1983}
S. Lang, \emph{Fundamentals of Diophantine geometry}, Springer, 1983,
Chapter 8, pp. 188--224.
\url{https://doi.org/10.1007/978-1-4757-1810-2}.

\bibitem{p2lbEGHP2006}
D. Eisenbud, M. Green, K. Hulek, and S. Popescu,
\emph{Small schemes and varieties of minimal degree}, American Journal
of Mathematics \textbf{128} (2006), no. 6, 1363--1389.
\url{https://doi.org/10.1353/ajm.2006.0043}.
\url{https://eisenbud.github.io/papers/pdfs/2006-002.pdf}.

\bibitem{p2lbHarrisTschinkel2000}
J. Harris and Y. Tschinkel, \emph{Rational points on quartics},
Duke Mathematical Journal \textbf{104} (2000), no. 3, 477--500.
\url{https://arxiv.org/abs/math/9809015}.
\url{https://cims.nyu.edu/~tschinke/papers/yuri/00quart/quart.pdf}.

\bibitem{p2lbBergEtAl2024}
J. Berg, C. Pagano, B. Poonen, M. Stoll, N. Triantafillou, B. Viray,
and I. Vogt, \emph{Brauer--Manin obstructions requiring arbitrarily many
Brauer classes}, Bulletin of the London Mathematical Society
\textbf{56} (2024), no. 5, 1587--1604; Theorem 5.1 and Section 6.1.
\url{https://doi.org/10.1112/blms.12999}.
\url{https://math.mit.edu/~poonen/papers/many_brauer.pdf}.

\bibitem{p2lbSerre1981}
J.-P. Serre, \emph{Quelques applications du th\'eor\`eme de densit\'e
de Chebotarev}, Publications Math\'ematiques de l'IH\'ES \textbf{54}
(1981), 123--201; Section 2.1, Theorem 1.
\url{https://numdam.org/item/PMIHES_1981__54__123_0.pdf}.

\bibitem{p2lbVoight2019}
J. Voight, \emph{Identifying the matrix ring: algorithms for quaternion
algebras and quadratic forms}, author's revised text, 21 May 2019,
Proposition 5.5; arXiv:1004.0994.
\url{https://jvoight.github.io/articles/quatalgs-051919.pdf}.

\bibitem{KoymansPagano2026}
P. Koymans and C. Pagano, \emph{Hilbert's tenth problem via additive
combinatorics}, Journal of the American Mathematical Society,
published electronically 10 August 2026.
\url{https://doi.org/10.1090/jams/1081}.
\url{https://arxiv.org/html/2412.01768v3}.

\bibitem{AlpogeBhargavaHoShnidman2026}
L. Alp\"oge, M. Bhargava, W. Ho, and A. Shnidman,
\emph{Rank stability in quadratic extensions and Hilbert's tenth
problem for the ring of integers of a number field}, Inventiones
Mathematicae \textbf{243} (2026), 1129--1139.
\url{https://doi.org/10.1007/s00222-025-01392-3}.
The cited Corollary 1.2 is that of the published version; the arXiv
version 1 states a weaker Diophantine-model conclusion.

\bibitem{Mazur1995}
B. Mazur, \emph{Speculations about the topology of rational points:
an up-date}, Ast\'erisque \textbf{228} (1995), 165--181;
Conjecture 3, p. 165.
\url{https://numdam.org/item/AST_1995__228__165_0/}.

\bibitem{Pan2022}
L. Pan, \emph{The Fontaine--Mazur conjecture in the residually
reducible case}, Journal of the American Mathematical Society
\textbf{35} (2022), 1031--1169.
\url{https://doi.org/10.1090/jams/991}.
\url{https://arxiv.org/abs/1901.07166}.

\bibitem{Yuan2024}
X. Yuan, \emph{Explicit Kodaira--Spencer map over Shimura curves},
Acta Mathematica Sinica, Chinese Series \textbf{67} (2024), no. 2,
227--249. \url{https://doi.org/10.12386/A20220154}.
\url{https://arxiv.org/abs/2205.11334}.

\end{thebibliography}
\end{document}
